% SHARD-ID: AQM-18-STALKS-AS-FIELD-GERMS
% SHARD-TITLE: Stalks as Field Germs
% SHARD-SUMMARY: Refines the phrase ``field at a point'' into stalk, local germ, fiber, and residue value.
% SHARD-SUMMARY: Applies the dictionary to \(\mathcal O(d)\otimes\mathbb F_3^2\) on \(\mathbf P^1_{\mathbb F_3}\).
% SHARD-SUMMARY: Tabulates every scalar residue used by the \(d=0,\ldots,4\) projective-line examples.
% SHARD-KEYWORDS: stalk, fiber, residue field, germ, projective line, arithmetic quantum field

\section{Stalks as Field Germs}
\label{sec:stalks-as-field-germs}

\subsection{Answer to the Terminology Question}

It is useful, but slightly dangerous, to say that a stalk is the ``field at the
point.''  The precise dictionary used here is:
\[
  \text{stalk}=\text{local field germ at the point},
  \qquad
  \text{fiber}=\text{actual value after residue reduction}.
\]
The source definition of a stalk as germs is
\path{references/algebraic_geometry/stacks_project_sheaves.tex}, lines
894--944.  The source Proj stalk formula is
\path{references/algebraic_geometry/stacks_project_constructions.tex}, lines
1091--1116 and 1182--1207.  The checked finite table is
\path{runs/2026-05-24-projective-line-sheaf-fields/data/projective_line_stalk_rows.csv}.

Let \(x\in X\) and let \(\mathcal E\) be a sheaf of
\(\mathcal O_X\)-modules.  The stalk \(\mathcal E_x\) consists of germs of
sections near \(x\).  If \(X\) is a scheme, \(\mathcal O_x\) is a local ring,
with maximal ideal \(\mathfrak m_x\), and residue field
\[
  \kappa(x)=\mathcal O_x/\mathfrak m_x.
\]
The fiber of \(\mathcal E\) at \(x\) is
\[
  \mathcal E(x)=\mathcal E_x\otimes_{\mathcal O_x}\kappa(x).
\]
Thus \(\mathcal E_x\) remembers infinitesimal or algebraic variation near
\(x\), while \(\mathcal E(x)\) is the value obtained after setting the local
coordinates equal to their value at \(x\).

\subsection{The Projective-Line Field Sheaf}

In the current example
\[
  X=\mathbf P^1_{\mathbb F_3},\qquad
  V=\mathbb F_3q\oplus\mathbb F_3p,\qquad
  \mathcal E_d=\mathcal O(d)\otimes_{\mathbb F_3}V.
\]
The rational points are ordered as
\[
  P_\infty=[1:0],\quad P_0=[0:1],\quad P_1=[1:1],\quad P_2=[2:1].
\]
For every rational point \(P\), the residue field is \(\kappa(P)=\mathbb F_3\).
Hence
\[
  \mathcal E_d(P)
  =
  \mathcal E_{d,P}\otimes_{\mathcal O_P}\kappa(P)
  \cong
  \mathbb F_3q\oplus\mathbb F_3p=V.
\]
This is the exact sense in which the arithmetic quantum field has a
point-value in \(V\).

The stalks themselves are larger local modules.  With
\[
  u=Y/X\quad\text{on }D_+(X),\qquad
  v=X/Y\quad\text{on }D_+(Y),
\]
and local frames \(e_X^{(d)}\), \(e_Y^{(d)}\), the four rational-point stalks
are:
\[
\begin{array}{c|c|c|c}
P & \mathcal O_P & \mathcal O(d)_P & \mathcal E_{d,P}\\
\hline
P_\infty & \mathbb F_3[u]_{(u)}
  & \mathbb F_3[u]_{(u)}e_X^{(d)}
  & \mathbb F_3[u]_{(u)}e_X^{(d)}\otimes V\\
P_0 & \mathbb F_3[v]_{(v)}
  & \mathbb F_3[v]_{(v)}e_Y^{(d)}
  & \mathbb F_3[v]_{(v)}e_Y^{(d)}\otimes V\\
P_1 & \mathbb F_3[v]_{(v-1)}
  & \mathbb F_3[v]_{(v-1)}e_Y^{(d)}
  & \mathbb F_3[v]_{(v-1)}e_Y^{(d)}\otimes V\\
P_2 & \mathbb F_3[v]_{(v-2)}
  & \mathbb F_3[v]_{(v-2)}e_Y^{(d)}
  & \mathbb F_3[v]_{(v-2)}e_Y^{(d)}\otimes V.
\end{array}
\]

\subsection{Germ-to-Value Maps}

Let
\[
  s_{d,i}=X^{d-i}Y^i,\qquad 0\leq i\leq d,
\]
and let \(w=aq+bp\in V\).  The vector-valued section is
\[
  \Phi_{d,i,w}=s_{d,i}\otimes w.
\]
Its stalk germ and fiber value at the four rational points are:
\[
\begin{array}{c|c|c}
P & \text{germ in }\mathcal E_{d,P} & \text{fiber value in }V\\
\hline
P_\infty
  & u^i e_X^{(d)}\otimes w
  & \begin{cases}w,&i=0,\\0,&i>0,\end{cases}\\
P_0
  & v^{d-i} e_Y^{(d)}\otimes w
  & \begin{cases}w,&i=d,\\0,&i<d,\end{cases}\\
P_1
  & v^{d-i} e_Y^{(d)}\otimes w
  & w\\
P_2
  & v^{d-i} e_Y^{(d)}\otimes w
  & 2^{d-i}w.
\end{array}
\]
The last line uses \(v=2=-1\) in \(\mathbb F_3\).  If \(w=(a,b)\) in the
ordered basis \((q,p)\), multiplication by the scalar \(\lambda\in\mathbb F_3\)
means
\[
  \lambda w=(\lambda a,\lambda b).
\]
These maps are not extra choices: they are the quotient maps
\(\mathcal E_{d,P}\to\mathcal E_d(P)\) induced by
\(\mathcal O_P\to\kappa(P)\).

\subsection{All Scalar Residues for Degrees 0 Through 4}

The following table gives the scalar \(\lambda_{d,i}(P)\) such that
\[
  \Phi_{d,i,w}(P)=\lambda_{d,i}(P)\,w.
\]
It is exactly the stalk-residue layer underlying the evaluation matrices in
Shard \(AQM\)-16.
\[
\begin{array}{c|c|c|c|c|c}
d & s_{d,i} & P_\infty & P_0 & P_1 & P_2\\
\hline
0 & 1       & 1 & 1 & 1 & 1\\
\hline
1 & X       & 1 & 0 & 1 & 2\\
1 & Y       & 0 & 1 & 1 & 1\\
\hline
2 & X^2     & 1 & 0 & 1 & 1\\
2 & XY      & 0 & 0 & 1 & 2\\
2 & Y^2     & 0 & 1 & 1 & 1\\
\hline
3 & X^3     & 1 & 0 & 1 & 2\\
3 & X^2Y    & 0 & 0 & 1 & 1\\
3 & XY^2    & 0 & 0 & 1 & 2\\
3 & Y^3     & 0 & 1 & 1 & 1\\
\hline
4 & X^4     & 1 & 0 & 1 & 1\\
4 & X^3Y    & 0 & 0 & 1 & 2\\
4 & X^2Y^2  & 0 & 0 & 1 & 1\\
4 & XY^3    & 0 & 0 & 1 & 2\\
4 & Y^4     & 0 & 1 & 1 & 1.
\end{array}
\]

\subsection{Concrete Reading of the Table}

For \(d=4\), the section \(X^3Y\otimes(q+2p)\) has:
\[
\begin{array}{c|c|c}
P & \text{germ} & \text{value}\\
\hline
P_\infty & u e_X^{(4)}\otimes(q+2p) & 0\\
P_0 & v^3 e_Y^{(4)}\otimes(q+2p) & 0\\
P_1 & v^3 e_Y^{(4)}\otimes(q+2p) & q+2p\\
P_2 & v^3 e_Y^{(4)}\otimes(q+2p) & 2q+p.
\end{array}
\]
The \(P_2\) value is \(2(q+2p)=2q+4p=2q+p\) in \(\mathbb F_3\).

Thus a stalk is the correct local object if one wants to discuss the field
near a point.  The finite Weyl label construction samples the residue/fiber
values, not the whole local stalk module.
